\documentclass[master,english]{include/tfgitic}

\usepackage{graphicx}
\usepackage{subcaption}
\usepackage{booktabs}
\usepackage{array}

\addbibresource{misc/thesis.bib}

\title{Estimating games of ordered preference}
\subtitle{Scalable preference-order estimation for Interactive Autonomous Systems}

\author[m]{Eric Roy-Almonacid}
\advisor[M]{Georgios Bakirtzis and Aleix Llusa Serra}
\school{Universitat Politecnica de Catalunya}
\degree{Master's degree in Machine Learning and Cybersecurity for Internet Connected Systems}
\location{Manresa (Barcelona), Spain}
\ects{18}
\topics{Autonomous Driving, Game-Theoretic Planning}

\begin{acknowledgments}
I would like to thank my advisor, Georgios Bakirtzis, for introducing me into
the research world and allowing me to attend a conference and a summer school.
I also want to thank Pau de las Heras Molins for allowing me to work together
in this exciting (but also intense) journey.
This thesis would've never seen the light without you both.

Moreover, this project wouldn't have seen the light if it wasn't for the
wonderful team that helped Pau and I submit our first paper, allowed us to
participate in their other projects, and gave feedback on our future works:
David Fridovich-Keil, Lasse Peters, Dong Ho Lee, and Georgios Bakirtzis.

I wish to thank too the board of the UPC-MERIT project and EPSEM's personnel to
partially fund my attendance to the conference and summer school. 
This work is partially supported by the E.U. Digital Europe Programme under grant
DIGITAL-101083531-MERIT.

Finally, I want to thank my family and friends for supporting my career and for
understanding my lack of social life these past months.
\end{acknowledgments}

\begin{resum}
TODO
\end{resum}

\begin{abstract}
Autonomous vehicles must balance ranked objectives, such as minimizing travel
time, ensuring safety, and coordinating with traffic.
Games of ordered preference effectively model these interactions but become
computationally intractable as the time horizon, number of players, or number
of preference levels increase.
Moreover, solving these games requires knowing the preferences of other agents.
% While receding horizon frameworks mitigate long-horizon intractability by solving sequential shorter games,
% often warm-started, they do not resolve the complexity growth inherent in existing methods for solving games of ordered preference.
This thesis introduces a solution strategy that avoids excessive complexity growth by approximating solutions using
lexicographic iterated best response in receding horizon, which uses past information to accelerate convergence.
Some methods are proposed to infer the preference order of other agents based on past observations.
The effectiveness of the proposed algorithms are demonstrated through simulated traffic scenarios, where approximate-optimal solutions for receding horizon games of ordered preference are computed, which converge towards generalized Nash equilibria.
A test environment using Duckietown's environment.
\end{abstract}

\begin{document}

\chapter{Introduction}

%% ITSC Introduction

% Complex agent decisions are often characterized by conflicting objectives.
% An autonomous vehicle, for example, must avoid collisions while also staying
% on the road, reaching a goal position,
% and obeying the speed limit.
% This problem is even
% harder when objectives of multiple agents conflict.
% Agents must determine which strategies are preferable,
% for example, going off-road
% to prevent a crash with another vehicle.
% Such decisions reflect an underlying comparative evaluation
% that ranks possible trajectories according
% to subjective preferences.
% We study the problem
% of comparative evaluations
% in transportation systems 
% from the lens of \emph{preference relations}
% that codify prioritized metrics; metrics
% that human drivers implicitly follow on the road.

% A preference is a ``total subjective comparative evaluation'' \cite{hausman:2011}.
% In the context
% of an autonomous vehicle, a total subjective comparative evaluation is a rule
% that, given two trajectories, decides which
% of the two is preferred from the perspective
% of the controller.
% For multiagent systems too,
% preferences model the multiobjective
% and often conflicting requirements
% that the system must adhere to
% for producing good behavior
% as subjective comparative evaluations (\cref{fig:scenario}).

% \begin{figure}[!t]
% \def\svgwidth{\linewidth}
% \subimport{figures/}{scenario.pdf_tex}
% \caption{Agents operate under shared constraints, but are incentivized by different preferences
% (adapted from Zanardi et al.~\protect\cite{zanardi:2022}, Mendes Filho et al.~\protect\cite{filho:2017}, and Lee et al.~\protect\cite{lee:2025}). On the road, situations arise where a vehicle has to coordinate with other vehicles and must relax a low-priority metric to preserve more important metrics, such as giving way to an ambulance.}
% \label{fig:scenario}
% %\vspace{-1em}
% \end{figure}

% One way to express the tension between prioritized metrics
% of multiple agents is to express them
% as lexicographic optimization problems
% in a game of ordered preference \cite{lee:2025}.
% Games of ordered preference are multi-player games
% with prohibitive computational cost.
% Solving these games with a receding horizon alleviates this cost
% through shorter time horizons and the introduction of feedback.
% However, the receding horizon setting does not address the rapid dimensional increase caused
% by adding players or preference levels.
% Iterated best response~(IBR) efficiently solves multi-player games
% by decomposing them into single-player problems,
% where each player optimizes assuming fixed opponent strategies.
% IBR convergence can be slow due
% to the need to explore large opponent decision spaces.
% We propose ``lexicographic IBR over time,''
% an extension of IBR that incorporates past information,
% available in receding horizon,
% to make predictions about future opponent decisions
% that accelerate convergence
% to generalized Nash equilibria
% for games of ordered preference.

% To study this area, this paper asks:
% \begin{enumerate}[itemsep=0pt,topsep=0pt,label=\textlangle ?\textrangle]
%     \item \emph{%
%         Can we efficiently compute approximate-optimal trajectories
%         for games of ordered preference?}\label{Q1}
% \end{enumerate}

% To make progress on this question, we
% \begin{enumerate}
% \item find that lexicographic IBR over time approximates optimal solutions
% for games of ordered preference; and
% \item design experiments demonstrating that games
% of ordered preference are an effective tool
% for analyzing interactive preferences in transportation systems.
% \end{enumerate}


% RAL

Conflicting objectives are present when having to make a decision. Autonomous agents need to make a choice of what objective should be prioritized. In urban driving scenarios, agents may prefer staying on the lane over staying under the speed limit to avoid hitting pedestrians, while the opposite may be preferred in the countryside.
Games of ordered preference model such scenarios in X way.

Non-cooperative games are often characterized by the lack of information
shared between different agents. In driving scenarios, goal destinations are rarely
broadcasted, enforcing an autonomous agent to make estimations on the
other agents' objectives.

While many works succeed in inferring some of the objectives, such as goal
position, there are no formal methods of estimating the preference order of other agents.

% \begin{figure}[!t]
% \def\svgwidth{\linewidth}
% \subimport{figures/}{scenario.pdf_tex}
% \caption{The order of the game objectives determines the behavior of an
% agent. To this end, estimating that order from partial state 
% observations allows the usage of games of ordered preference in
% environments where there is little to no communications between agents,
% such as traffic scenarios with humans involved.}
% \label{fig:scenario}
% %\vspace{-1em}
% \end{figure}


% To study this area, this paper asks:
% \begin{enumerate}[itemsep=0pt,topsep=0pt,label=\textlangle ?\textrangle]
%     \item \emph{%
%         Can we efficiently infer other agents' preference relations in
%         games of ordered preference?}\label{Q1}
% \end{enumerate}

% To make progress on this question, we
% 1. design a method that estimates preference relations in games of ordered
%    preference; and
% 2. test the previous tool on small-scale robots (cref:fig:scenario), demonstrating
%    the utility of preference estimation.

\chapter{Research Context and Thesis Scope}

\section{Context}

This thesis is positioned in interactive autonomous systems, where agents must
make decisions under competing objectives and limited information. The central
technical lens is games of ordered preference (GOOPs), which encode strict
objective priorities instead of weighted trade-offs \cite{lee:2025,zanardi:2022}.

The practical motivation is to bridge three levels that are usually separated:

\begin{enumerate}
    \item formal game-theoretic planning,
    \item deployable robotics software,
    \item inverse inference of other agents' hidden preferences.
\end{enumerate}

\section{Thesis question}

The unifying question is:

\begin{quote}
How can ordered-preference game formulations be made practical for real
interactive systems when computation is constrained and other agents do not
reveal their preference structure?
\end{quote}

\section{Subproblems addressed}

This thesis decomposes that question into three subproblems.

\begin{description}
    \item[SP1: scalable forward solving] Efficient approximate solving of GOOPs
    in receding horizon.
    \item[SP2: embodied integration] Execution of the solver stack in
    Duckietown through a robust manager--solver--adapter architecture.
    \item[SP3: inverse preference estimation] Estimation of objective ordering
    from observed behavior in GOOP-like settings.
\end{description}

\section{Artifacts used in this thesis}

The manuscript is built around four concrete artifacts:

\begin{enumerate}
    \item the ITSC paper on approximate GOOP solving \cite{delasheras:2025},
    \item the \texttt{patos} integration repository,
    \item the \texttt{InverseGoopSolver} repository,
    \item the ongoing RA-L manuscript draft on preference estimation.
\end{enumerate}

\section{What this draft delivers}

At the current stage, this thesis draft provides:

\begin{itemize}
    \item full technical foundations and notation,
    \item complete software architecture description,
    \item implemented inverse-estimation workflow baseline,
    \item explicit experimental and writing roadmap to final submission.
\end{itemize}

Because parts of SP2 and SP3 are still in progress, this draft emphasizes
methodological clarity and reproducible structure rather than final
performance claims.

\begin{figure}[tbp]
    \centering
    \fbox{\parbox{0.92\linewidth}{
    \textbf{Placeholder figure: Thesis work map.}
    A map connecting SP1/SP2/SP3 with repositories, experiment scripts, and
    chapter outputs.
    }}
    \caption{Scope and structure of the MSc thesis work.}
    \label{fig:scope_map}
\end{figure}

\section{Contribution framing}

The contribution of this thesis is not a single isolated algorithm. Instead, it
is an integrated pipeline:

\begin{center}
\emph{approximate forward solving} $\rightarrow$
\emph{robotics integration} $\rightarrow$
\emph{inverse preference estimation}.
\end{center}

This framing is the key distinction from the earlier bachelor thesis and from a
paper-only manuscript.

\chapter{Mathematical Foundations}

\section{Trajectory-game model}

Let $i \in [N]$ index agents. At time $t$, each agent state is
$x_t^i \in \mathcal{X}^i$, control is $u_t^i \in \mathcal{U}^i$, and decision
variable is $z_t^i = [x_t^i, u_t^i]$. Over horizon $T$, trajectory
$\mathbf{z}^i = \{z_0^i,\dots,z_{T-1}^i\}$.

A constrained trajectory game for agent $i$ is:
\begin{equation}
\begin{aligned}
\min_{\mathbf{z}^i} \quad & J^i(\mathbf{z}^i,\mathbf{z}^{-i};\theta^i) \\
\text{s.t.}\quad & x_{t+1}^i = f^i(x_t^i,u_t^i), \\
& x_0^i = \hat{x}_0^i, \\
& g^i(\mathbf{z})=0,\; h^i(\mathbf{z})\ge0, \\
& g^s(\mathbf{z})=0,\; h^s(\mathbf{z})\ge0.
\end{aligned}
\label{eq:traj_game_foundation}
\end{equation}

\section{Ordered preferences and lexicographic optimization}

Each agent has an ordered objective list
$\{J_1^i, J_2^i, \dots, J_{K^i}^i\}$, where lower index means higher priority.
GOOP solutions enforce strict precedence through nested optimization
\cite{lee:2025}:
\begin{equation}
\begin{aligned}
\min_{\mathbf{z}_{K^i}^i} \;& J_{K^i}^i(\mathbf{z}_{K^i}^i,\mathbf{z}^{-i}) \\
\text{s.t.}\;& \mathbf{z}_{K^i}^i \in \arg\min_{\mathbf{z}_{K^i-1}^i} J_{K^i-1}^i(\cdot), \\
& \vdots \\
& \mathbf{z}_2^i \in \arg\min_{\mathbf{z}_1^i} J_1^i(\mathbf{z}_1^i,\mathbf{z}^{-i}),
\end{aligned}
\label{eq:goop_nested_foundation}
\end{equation}
plus private and shared feasibility constraints.

\section{Receding-horizon execution}

Long-horizon solves are split into repeated short-horizon solves of length $T$
with execution chunk $T_l \le T$ \cite{christophersen:2007,hall:2022}. At each
step:

\begin{enumerate}
    \item measure current joint state,
    \item solve GOOP over window $T$,
    \item execute first $T_l$ controls,
    \item shift and repeat.
\end{enumerate}

\section{Lexicographic IBR over time}

The baseline method used in this thesis is lexicographic IBR over time
\cite{delasheras:2025}:

\begin{itemize}
    \item decompose multi-agent GOOP into per-agent lexicographic best
    responses,
    \item iterate best responses for a bounded number of rounds $L$,
    \item warm start each step with shifted trajectories from previous solves.
\end{itemize}

This introduces a runtime-quality trade-off controlled by $L$ and makes larger
receding-horizon problems tractable in practice.

\section{Bridge to inverse estimation}

The same model supports inverse questions: if objective primitives are known but
ordering is unknown, estimate ordering by matching predicted and observed
behavior over an identification horizon. This bridge is central to the
\texttt{InverseGoopSolver} extension.

\begin{figure}[tbp]
    \centering
    \fbox{\parbox{0.92\linewidth}{
    \textbf{Placeholder figure: Forward/inverse shared model.}
    Diagram showing that forward solving and inverse estimation share dynamics,
    constraints, and objective primitives, differing mainly in which variables
    are optimized versus inferred.
    }}
    \caption{Shared mathematical core across thesis work packages.}
    \label{fig:shared_core}
\end{figure}

\chapter{Published Baseline: Approximate GOOP Solving}

\section{Chapter role}

This chapter documents the completed and published foundation of the thesis:
\emph{Approximate solutions to games of ordered preference}
\cite{delasheras:2025}. The rest of the manuscript extends this baseline rather
than replacing it.

\section{Core idea}

The published method combines receding horizon with lexicographic IBR over time
for approximate GOOP solving. The motivation is to avoid the cost of repeatedly
solving a tightly coupled multi-agent lexicographic program at each control
cycle.

\section{Algorithmic structure}

For each receding-horizon step:

\begin{enumerate}
    \item initialize trajectory predictions for all agents from shifted previous
    solutions,
    \item run bounded IBR rounds,
    \item in each round, each agent solves its lexicographic best-response
    subproblem against fixed predictions,
    \item aggregate resulting controls and execute the next segment.
\end{enumerate}

This process keeps preference semantics explicit while reducing per-step
computational burden.

\section{Implementation and tooling}

The implementation is in Julia \cite{julia} and uses the GOOP ecosystem around
TrajectoryGamesBase.jl \cite{peters:2024a}, ParametricMCPs.jl
\cite{peters:2025}, and PATH-backed complementarity solving
\cite{dirkse:1995}. The released package is LEXIBROverTime.jl \cite{code}.

\section{Representative results}

The paper reports substantial runtime gains compared to a tightly coupled
baseline while preserving approximate solution quality over practical iteration
budgets.

\begin{table}[tbp]
\centering
\caption{Representative values reported in \cite{delasheras:2025}.}
\label{tab:published_runtime}
\begin{tabular}{@{}lrrrr@{}}
\toprule
 & \multicolumn{2}{c}{$K=2$} & \multicolumn{2}{c}{$K=3$} \\
\cmidrule(r){2-3}\cmidrule(l){4-5}
Method & Solve time [s] & $L_1$ distance & Solve time [s] & $L_1$ distance \\
\midrule
Baseline & 44.51 & -- & 2676.00 & -- \\
$L=1$ & 0.36 & $3.28\times10^{-4}$ & 0.32 & $5.45\times10^{-4}$ \\
$L=2$ & 0.59 & $1.22\times10^{-5}$ & 0.63 & $6.56\times10^{-5}$ \\
$L=5$ & 1.39 & $1.78\times10^{-6}$ & 1.20 & $4.96\times10^{-5}$ \\
\bottomrule
\end{tabular}
\end{table}

\section{What is reused in this thesis}

This baseline contributes the reusable components adopted in later chapters:

\begin{itemize}
    \item receding-horizon execution interface,
    \item trajectory warm-start logic,
    \item lexicographic objective structure,
    \item solver hyperparameterization and diagnostics.
\end{itemize}

\section{Open points left by the paper}

The paper leaves two practical gaps addressed by the MSc thesis:

\begin{enumerate}
    \item deployment-level integration with physical/simulated robots,
    \item inverse estimation of hidden preference structure.
\end{enumerate}

These are exactly the focus of the following chapters.

\chapter{Duckietown System Integration (\texttt{patos})}

\section{Objective}

This chapter covers the engineering extension that brings the planning stack
into an executable Duckietown workflow. The goal is robust closed-loop operation,
not only offline trajectory generation.

\section{Runtime architecture}

The runtime is organized into three coordinated modules:

\begin{enumerate}
    \item manager node (state aggregation and acceleration publication),
    \item Julia TCP bridge (receding-horizon solver service),
    \item per-robot adapter node (acceleration-to-command conversion).
\end{enumerate}

\begin{figure}[tbp]
    \centering
    \fbox{\parbox{0.92\linewidth}{
    \textbf{Placeholder figure: \texttt{patos} runtime architecture.}
    Include ROS topics and TCP request/response path between manager and Julia
    bridge.
    }}
    \caption{Runtime architecture of the \texttt{patos} integration.}
    \label{fig:patos_runtime_arch}
\end{figure}

\section{Manager behavior}

The manager supports multi-robot operation and can run in:

\begin{itemize}
    \item trajectory-sequence mode (CSV-driven),
    \item Julia-solver mode (online solver queries).
\end{itemize}

In solver mode, it sends per-robot $(x,y,v_x,v_y,\theta)$ snapshots and
receives $(a_x,a_y)$ commands. Error handling is explicit: on solver timeout or
failure, the manager can publish zero acceleration or fallback controls depending
on configuration.

\section{Adapter behavior}

Each adapter converts incoming acceleration vectors into bounded
velocity/yaw-rate commands using an internal velocity-vector controller with
state decay and heading correction. The adapter also publishes telemetry,
including model state, command state, and optional encoder/wheel-command
feedback.

This makes the controller chain observable and debuggable for both simulation
and physical runs.

\section{Julia bridge details}

The bridge maintains persistent context for receding-horizon progression and
provides simple commands (e.g., \texttt{PING}, \texttt{STEP}, \texttt{RESET}).
Internally it builds and solves a constrained multi-player planning problem with
road, speed, and collision constraints plus goal-seeking objectives.

\section{Deployment workflow}

The repository includes scripts/UI for:

\begin{itemize}
    \item container build and launch,
    \item code transfer to one or multiple robots,
    \item starting/stopping manager and adapters,
    \item telemetry inspection and map-level visualization.
\end{itemize}

\section{Validation path}

Validation is layered:

\begin{enumerate}
    \item unit-level controller tests,
    \item local simulation checks,
    \item ROS container smoke tests,
    \item real robot runs.
\end{enumerate}

This sequence reduces integration risk and supports repeatable debugging.

\section{Current maturity}

The system is functional and can run end-to-end. Remaining work for final thesis
results is focused on broader scenario coverage and quantitative robustness
metrics.

\section{Positioning}

This work extends simulation-first GOOP literature by providing a concrete
systems bridge to an embodied platform \cite{duckietown:2017,delasheras:2025}.

\chapter{Inverse Preference Estimation Extension}

\section{Motivation}

Forward GOOP planning assumes that objective hierarchies are known. In realistic
traffic-like interaction, this is not true: other agents typically do not expose
their preference ordering. This chapter introduces the current estimation
extension built in \texttt{InverseGoopSolver}.

\section{Repository workflow}

The extension is organized into three executable programs:

\begin{description}
    \item[Program 1] centralized forward baseline,
    \item[Program 2] decentralized per-agent solving with prediction beliefs,
    \item[Program 3] ego preference-order identification by candidate testing.
\end{description}

\section{Program 3 estimation logic}

Program 3 currently estimates ordering by permutation search over a candidate
set for the ego agent:

\begin{enumerate}
    \item generate an oracle rollout with known ordering,
    \item test candidate orderings over an identification horizon,
    \item compute mismatch score between predicted and observed controls,
    \item select minimum-score ordering,
    \item run decentralized simulation with inferred order.
\end{enumerate}

Current score:
\begin{equation}
S(\pi) = \sum_{t=1}^{T_{\mathrm{id}}}
\|u^{\mathrm{pred}}_{\pi,t} - u^{\mathrm{obs}}_t\|_2^2.
\label{eq:score_inverse_chapter}
\end{equation}

\begin{figure}[tbp]
    \centering
    \fbox{\parbox{0.92\linewidth}{
    \textbf{Placeholder figure: Inverse estimation pipeline.}
    Show oracle rollout, candidate-order loop, score computation
    (Eq.~\ref{eq:score_inverse_chapter}), best-order selection, and
    decentralized rerun.
    }}
    \caption{Current inverse preference-order estimation workflow.}
    \label{fig:inverse_pipeline}
\end{figure}

\section{Relation to ongoing manuscript}

The ongoing RA-L manuscript direction uses this implementation as an executable
baseline to study preference-relation inference under partial observations. The
thesis keeps this work explicit as \emph{in-progress research}, not as finished
final results.

\section{Current limitations}

Main limitations at this stage:

\begin{itemize}
    \item factorial candidate growth with objective count,
    \item limited robustness analysis under noisy/partial observations,
    \item partial integration with embodied Duckietown experiments.
\end{itemize}

\section{Planned upgrades}

Planned improvements for final thesis stage include candidate pruning,
uncertainty-aware scoring, and broader benchmark coverage
\cite{peters:2023,li:2023,liu:2024}.

\chapter{Experimental Methodology}

\section{Evaluation philosophy}

The thesis evaluation strategy is incremental:

\begin{enumerate}
    \item validate solver behavior in controlled simulation,
    \item validate software integration under runtime constraints,
    \item validate inverse estimation quality under increasing uncertainty.
\end{enumerate}

This staged methodology prevents mixing algorithmic and systems failures in a
single experiment.

\section{Evaluation dimensions}

Three dimensions are measured across work packages:

\begin{itemize}
    \item \textbf{efficiency}: solve time and runtime overhead,
    \item \textbf{behavior}: safety/goal/lane metrics and qualitative
    trajectory signatures,
    \item \textbf{estimation}: preference-order identification quality and
    downstream prediction impact.
\end{itemize}

\section{Forward-solver benchmarks}

For forward solving, the baseline metrics from \cite{delasheras:2025} are reused
and extended:

\begin{itemize}
    \item solve time per receding-horizon step,
    \item compile/setup overhead,
    \item convergence or stabilization behavior versus IBR rounds.
\end{itemize}

\section{Integration benchmarks}

For \texttt{patos}, planned metrics include:

\begin{itemize}
    \item manager cycle consistency,
    \item solver timeout/fallback rate,
    \item command tracking and safety-clearance indicators,
    \item behavior under multi-robot load.
\end{itemize}

\section{Inverse-estimation benchmarks}

For \texttt{InverseGoopSolver}, planned metrics include:

\begin{itemize}
    \item exact-match rate of inferred order,
    \item rank-distance metrics between inferred and true order,
    \item trajectory mismatch before/after estimation,
    \item robustness against observation noise and partial state access.
\end{itemize}

\begin{figure}[tbp]
    \centering
    \fbox{\parbox{0.92\linewidth}{
    \textbf{Placeholder figure: Metric dashboard layout.}
    Panel view with efficiency, safety/behavior, and estimation metrics to be
    reported uniformly across experiments.
    }}
    \caption{Planned unified metric dashboard for final results chapters.}
    \label{fig:metric_dashboard}
\end{figure}

\section{Reproducibility protocol}

To keep comparisons reliable, each final experiment set will record:

\begin{itemize}
    \item exact configuration parameters,
    \item software commit references,
    \item random seeds,
    \item generated plots and serialized summary artifacts.
\end{itemize}

This protocol will be used for both simulation runs and embodied Duckietown
experiments.

\chapter{Current Results, Gaps, and Risks}

\section{Status snapshot}

At this draft stage, the thesis has:

\begin{itemize}
    \item completed and published forward-solving baseline,
    \item functional Duckietown integration prototype,
    \item functional inverse-estimation prototype pipeline.
\end{itemize}

\section{What is already supported by evidence}

Evidence is currently strongest for:

\begin{enumerate}
    \item scalability gains of approximate GOOP solving
    \cite{delasheras:2025},
    \item end-to-end feasibility of the manager--bridge--adapter architecture,
    \item operational preference-order inference workflow.
\end{enumerate}

\section{What is still missing for final claims}

The following pieces are still needed for final thesis-level conclusions:

\begin{itemize}
    \item broader quantitative integration benchmarks,
    \item larger and noisier inverse-estimation benchmark sets,
    \item tighter cross-validation between simulation and embodied runs.
\end{itemize}

\section{Main technical risks}

\begin{enumerate}
    \item inverse search scalability with larger objective sets,
    \item runtime sensitivity to communication delays,
    \item scenario overfitting if benchmark diversity remains limited.
\end{enumerate}

\section{Risk mitigation}

Planned mitigation actions:

\begin{itemize}
    \item prioritize one primary benchmark protocol before adding variants,
    \item enforce fixed configuration templates for comparison,
    \item add stress-test suites early to detect integration bottlenecks.
\end{itemize}

\begin{figure}[tbp]
    \centering
    \fbox{\parbox{0.92\linewidth}{
    \textbf{Placeholder figure: Progress tracker.}
    A progress bar chart for SP1/SP2/SP3 with completed items, ongoing items,
    and remaining validation tasks.
    }}
    \caption{Current maturity and remaining work by thesis subproblem.}
    \label{fig:progress_tracker}
\end{figure}

\chapter{Current Status and Preliminary Discussion}

\section{Purpose of this chapter}

Because part of the thesis work is still ongoing, this chapter intentionally
provides a \emph{status-oriented discussion} instead of a final conclusion
chapter. The goal is to make progress explicit, identify current evidence, and
highlight what remains to reach thesis completion.

\section{Status by work package}

Table~\ref{tab:status_matrix} summarizes the current maturity of each work
package.

\begin{table}[tbp]
\centering
\caption{Current status matrix (draft stage).}
\label{tab:status_matrix}
\begin{tabular}{@{}p{0.22\linewidth}p{0.2\linewidth}p{0.5\linewidth}@{}}
\toprule
Work package & Status & Evidence available in this draft \\
\midrule
WP1: Approximate GOOP solving (paper) & Completed & Formal method, implementation, and published-style quantitative results \cite{delasheras:2025}. \\
WP2: Duckietown integration (\texttt{patos}) & Functional prototype & End-to-end software architecture, manager/adapter/bridge loop, deployment scripts, and test workflow documentation. \\
WP3: Inverse estimation (\texttt{InverseGoopSolver}) & Exploratory prototype & Centralized/decentralized runners and candidate-based preference-order inference workflow with reproducible outputs. \\
\bottomrule
\end{tabular}
\end{table}

\section{What is already demonstrated}

At this stage, the thesis already supports the following claims with concrete
artifacts:

\begin{enumerate}
    \item Approximate lexicographic IBR over time can drastically reduce solve
    time versus tightly coupled baseline solves in representative GOOP scenarios.
    \item The same planning concepts can be integrated into a practical robotics
    stack with clear interfaces and runtime safety fallbacks.
    \item Preference-order estimation can be operationalized as a reproducible
    inverse workflow on top of the same planning backbone.
\end{enumerate}

These claims are currently strongest for algorithmic feasibility and software
architecture. They are weaker for broad statistical generalization, which
requires additional experiments.

\section{Main limitations of the current draft}

The current draft has four important limitations:

\begin{itemize}
    \item \textbf{Incomplete WP2 quantitative evaluation.} Runtime robustness and
    closed-loop quality metrics in broader multi-robot conditions are still being
    collected.
    \item \textbf{Incomplete WP3 statistical evaluation.} Preference-order
    estimation is implemented but not yet benchmarked at the scale needed for
    final claims.
    \item \textbf{Combinatorial growth in inverse search.} Full permutation
    testing does not scale to large objective sets.
    \item \textbf{Partial integration of inverse methods with embodied tests.}
    The full bridge from inverse estimation to real-robot execution remains
    ongoing.
\end{itemize}

\section{Preliminary interpretation}

Even with these limitations, the current evidence suggests that the thesis
pipeline is coherent:

\begin{itemize}
    \item the forward method is fast enough to support repeated replanning,
    \item the software integration makes this method executable in a realistic
    robotics loop,
    \item inverse estimation can reuse the same formal objects and code
    abstractions, reducing conceptual and implementation fragmentation.
\end{itemize}

This coherence is important: it indicates that final thesis value is likely to
come from \emph{integration quality} as much as from any single algorithmic
result.

\section{Risks to final completion}

The main risks between this draft and the final thesis submission are:

\begin{enumerate}
    \item delays in collecting stable WP2 multi-robot metrics,
    \item insufficient statistical depth for WP3 evaluation,
    \item scope creep in inverse-model variants before consolidating a primary
    benchmark protocol.
\end{enumerate}

Mitigation actions are listed in the roadmap chapter.

\section{Takeaway}

This draft already establishes a strong technical baseline and integrated
architecture. The final thesis effort is now concentrated on completion-quality
evidence: broader experiments, consolidated metrics, and final synthesis across
forward solving, deployment, and inverse estimation.

\chapter{Roadmap to Final Submission}

\section{Completion plan}

From this draft to the final thesis version, the roadmap is:

\begin{enumerate}
    \item complete WP2 quantitative integration benchmarks,
    \item complete WP3 estimation robustness benchmarks,
    \item consolidate cross-package analysis and final conclusions,
    \item replace all placeholders with final figures and tables.
\end{enumerate}

\section{Planned milestones}

\begin{table}[tbp]
\centering
\caption{Planned finalization milestones.}
\label{tab:final_milestones}
\begin{tabular}{@{}p{0.22\linewidth}p{0.2\linewidth}p{0.5\linewidth}@{}}
\toprule
Window & Focus & Deliverable \\
\midrule
March--April 2026 & Integration metrics & Stable runtime and safety metrics for \texttt{patos}. \\
April--May 2026 & Estimation benchmarks & Consolidated inverse-order inference results under uncertainty. \\
May 2026 & Thesis synthesis & Final discussion and conclusion chapter updates. \\
Final revision & Manuscript polish & Final figure set, bibliography cleanup, defense-ready PDF. \\
\bottomrule
\end{tabular}
\end{table}

\section{Beyond-thesis extensions}

Potential post-thesis extensions:

\begin{itemize}
    \item scalable order-search methods,
    \item online preference-belief updates in receding horizon,
    \item richer mixed-traffic and intersection scenarios,
    \item stronger sim-to-real transfer evaluations.
\end{itemize}

\section{Expected final thesis claim}

The target final claim of this MSc thesis is:

\begin{quote}
Ordered-preference game models become practically useful when efficient
approximate solving, systems integration, and inverse preference estimation are
developed as one coherent stack.
\end{quote}


\printbibliography

\end{document}
